\documentclass[../DoAn.tex]{subfiles}
\begin{document}

Phụ lục này bổ sung đặc tả chi tiết cho bốn use case quan trọng nhưng chưa được đặc tả đầy đủ trong phần nội dung chính (Chương \ref{chapter:Related_works} mới đặc tả sáu use case trọng yếu UC-01 đến UC-06). Bốn use case đặc tả thêm dưới đây bao gồm Quản lý đơn vị (UC-07), Sao lưu dữ liệu (UC-08), Quản lý tài khoản (UC-09) và Xem nhật ký hệ thống (UC-10).

\section{Đặc tả use case ``Quản lý đơn vị (CQĐV / ĐVTT)''}
\label{appendix:uc07}

\begin{table}[H]
\centering
\caption{Đặc tả use case UC-07 Quản lý đơn vị}
\label{tab:uc07}
\begin{tabular}{|>{\raggedright\arraybackslash}p{3cm}|p{11.5cm}|}
\hline
\textbf{Mục} & \textbf{Nội dung} \\ \hline
Tên & Quản lý đơn vị (CQĐV / ĐVTT) \\ \hline
ID & UC-07 \\ \hline
Tác nhân & SuperAdmin (cấu hình hạ tầng), Admin (Phòng Chính trị quản trị nghiệp vụ) \\ \hline
Tiền điều kiện & Đã đăng nhập với vai trò có quyền chỉnh sửa cây tổ chức. \\ \hline
Luồng chính &
(1) Thêm cơ quan đơn vị (CQĐV) cấp trên: chọn ``Thêm CQĐV'', nhập \texttt{ma\_don\_vi}, \texttt{ten\_don\_vi}; hệ thống xác thực \texttt{ma\_don\_vi} duy nhất rồi lưu vào bảng \texttt{CoQuanDonVi}. (2) Thêm đơn vị trực thuộc (ĐVTT): chọn CQĐV cha, nhập thông tin ĐVTT; hệ thống lưu vào \texttt{DonViTrucThuoc} với khoá ngoại \texttt{co\_quan\_don\_vi\_id}. (3) Quản lý chức vụ trong CQĐV/ĐVTT: thêm chức vụ với \texttt{ten\_chuc\_vu}, \texttt{is\_manager}, \texttt{he\_so\_chuc\_vu}; hệ thống lưu kèm ràng buộc duy nhất theo cặp đơn vị -- tên chức vụ. (4) Xem cây đơn vị kèm bộ đếm số quân nhân: hệ thống trả về cấu trúc cây với \texttt{so\_luong} đã được tự động cập nhật khi thêm/sửa/xoá quân nhân. \\ \hline
Luồng thay thế &
(1a) \texttt{ma\_don\_vi} đã tồn tại $\rightarrow$ 400 ``Mã đơn vị đã tồn tại''. (2a) CQĐV cha không tồn tại $\rightarrow$ 404. (Xoá CQĐV) Hệ thống cascade xoá tất cả ĐVTT, ChucVu, QuanNhan thuộc CQĐV theo \texttt{onDelete: Cascade} ở schema. \\ \hline
Hậu điều kiện & Cấu trúc cây đơn vị được cập nhật; \texttt{so\_luong} chính xác; các quân nhân và chức vụ được liên kết đúng. \\ \hline
\end{tabular}
\end{table}

\section{Đặc tả use case ``Sao lưu dữ liệu''}
\label{appendix:uc08}

\begin{table}[H]
\centering
\caption{Đặc tả use case UC-08 Sao lưu dữ liệu}
\label{tab:uc08}
\begin{tabular}{|>{\raggedright\arraybackslash}p{3cm}|p{11.5cm}|}
\hline
\textbf{Mục} & \textbf{Nội dung} \\ \hline
Tên & Sao lưu dữ liệu định kỳ và thủ công \\ \hline
ID & UC-08 \\ \hline
Tác nhân & SuperAdmin (qua DevZone), Cron task tự động \\ \hline
Tiền điều kiện & Tính năng backup được bật (\texttt{SystemSetting.backup\_enabled = true}); SuperAdmin đã xác thực mật khẩu DevZone qua \texttt{POST /api/dev-zone/auth}. \\ \hline
Luồng chính &
(1) Cron job tự động: \texttt{node-cron} in-process khởi chạy theo \texttt{cron\_schedule} (mặc định \texttt{'0 1 1 * *'} --- 01:00 ngày 1 mỗi tháng); gọi \texttt{backup.service.ts.createBackup(\{ type: 'scheduled' \})}; đọc 21 bảng dữ liệu qua \texttt{Promise.all}; build chuỗi SQL \texttt{INSERT INTO ... VALUES (...)} cho từng row, escape JSON/string qua \texttt{quoteValue()}; \texttt{fs.writeFileSync(filePath, sqlLines.join('\textbackslash n'))} vào thư mục \texttt{backups/}; cập nhật \texttt{SystemSetting.backup\_last\_run = now ISO}; ghi \texttt{SystemLog} \texttt{resource = 'backup'} (chỉ SuperAdmin xem được); xoá file backup cũ hơn \texttt{backup\_retention\_days} (mặc định 15). (2) SuperAdmin trigger thủ công qua \texttt{POST /api/dev-zone/backup/trigger}, gọi \texttt{createBackup(\{ type: 'manual' \})} --- flow giống Cron. (3) Xem trạng thái: \texttt{GET /api/dev-zone/backup/status} trả về \texttt{\{ enabled, schedule, last\_run, next\_run, retention\_days, recent\_logs \}}. (4) Đổi lịch / bật-tắt: \texttt{PUT /api/dev-zone/backup/schedule} body \texttt{\{ enabled, schedule, retention\_days \}}; xác thực cú pháp cron, cập nhật \texttt{SystemSetting}, đăng ký lại \texttt{node-cron}. (5) Cleanup thủ công: \texttt{POST /api/dev-zone/backup/cleanup}, quét \texttt{backups/}, xoá file quá hạn. \\ \hline
Luồng thay thế &
(1a) \texttt{backup\_enabled = false} $\rightarrow$ cron skip silently, không sinh file. (1.4a) I/O fail (đầy ổ cứng / không có quyền) $\rightarrow$ ghi \texttt{SystemLog} \texttt{action = BACKUP\_FAILED}; không gửi notification (chỉ SuperAdmin xem được). (4a) Cú pháp cron không hợp lệ $\rightarrow$ trả 400, giữ schedule cũ. \\ \hline
Hậu điều kiện & File \texttt{.sql} mới được lưu trong thư mục \texttt{backups/} của BE-QLKT. Restore phải thủ công qua \texttt{psql -d qlkt < backup.sql}. \\ \hline
\end{tabular}
\end{table}

\section{Đặc tả use case ``Quản lý tài khoản''}
\label{appendix:uc09}

\begin{table}[H]
\centering
\caption{Đặc tả use case UC-09 Quản lý tài khoản}
\label{tab:uc09}
\begin{tabular}{|>{\raggedright\arraybackslash}p{3cm}|p{11.5cm}|}
\hline
\textbf{Mục} & \textbf{Nội dung} \\ \hline
Tên & Quản lý tài khoản người dùng \\ \hline
ID & UC-09 \\ \hline
Tác nhân & SuperAdmin, Admin (Phòng Chính trị) \\ \hline
Tiền điều kiện & Đã đăng nhập với vai trò SuperAdmin hoặc Admin; route \texttt{/api/accounts} dùng middleware \texttt{requireAdmin}. \\ \hline
Luồng chính &
(1) Tạo tài khoản mới: chọn quân nhân, nhập \texttt{username} và \texttt{role}; hệ thống sinh password mặc định, băm bằng bcrypt, lưu \texttt{TaiKhoan} với khoá ngoại \texttt{quan\_nhan\_id}; trả password gốc một lần để người tạo chuyển cho người dùng. (2) Cập nhật role hoặc khoá tài khoản: sửa \texttt{role} hoặc bật/tắt \texttt{active}; \texttt{UPDATE TaiKhoan}, ghi log. (3) Reset mật khẩu: hệ thống sinh password mới, băm, lưu, trả về một lần. (4) Xoá tài khoản: xác nhận; \texttt{DELETE TaiKhoan}, các log liên quan giữ lại với \texttt{SetNull}. \\ \hline
Luồng thay thế &
(1a) Username đã tồn tại $\rightarrow$ 400 ``Username đã được sử dụng''. (1b) Quân nhân đã có tài khoản (ràng buộc 1-1) $\rightarrow$ 400. \\ \hline
Hậu điều kiện & Tài khoản được tạo/cập nhật/xoá trong cơ sở dữ liệu; người dùng có thể đăng nhập với tài khoản mới. \\ \hline
\end{tabular}
\end{table}

\section{Đặc tả use case ``Xem nhật ký hệ thống''}
\label{appendix:uc10}

\begin{table}[H]
\centering
\caption{Đặc tả use case UC-10 Xem nhật ký hệ thống}
\label{tab:uc10}
\begin{tabular}{|>{\raggedright\arraybackslash}p{3cm}|p{11.5cm}|}
\hline
\textbf{Mục} & \textbf{Nội dung} \\ \hline
Tên & Xem nhật ký hệ thống \\ \hline
ID & UC-10 \\ \hline
Tác nhân & SuperAdmin, Admin, Manager \\ \hline
Tiền điều kiện & Đã đăng nhập với vai trò SuperAdmin, Admin hoặc Manager; route \texttt{/api/system-logs} dùng \texttt{requireManager}. Log có \texttt{resource = 'backup'} chỉ SuperAdmin xem được --- Admin và Manager bị filter trong \texttt{systemLogs.service.ts}. \\ \hline
Luồng chính &
(1) Truy cập trang ``Nhật ký hệ thống''; hệ thống hiển thị bảng log với cột: thời gian, người thực hiện, vai trò, action, resource, mô tả. (2) Áp dụng bộ lọc theo resource, action, người thực hiện, khoảng thời gian. (3) Hệ thống gọi \texttt{systemLogsService.getLogs(userRole, filters)} --- filter \texttt{resource = 'backup'} nếu role không phải SuperAdmin; Manager còn bị giới hạn theo phạm vi đơn vị qua \texttt{getManagerAccountIds}. (4) Trả về danh sách log đã phân trang. (5) Click vào một log để xem chi tiết \texttt{payload} (before/after JSON) trong modal. \\ \hline
Luồng thay thế &
(3a) Admin hoặc Manager cố xem log \texttt{resource = 'backup'} $\rightarrow$ tự động bị filter trong service, không thấy bản ghi nào. (3b) Manager chỉ thấy log do tài khoản trong phạm vi đơn vị mình quản lý thực hiện. \\ \hline
Hậu điều kiện & Người dùng nắm được lịch sử hành động trong hệ thống, phục vụ công tác kiểm tra và truy vết. \\ \hline
\end{tabular}
\end{table}

\end{document}
